\section{Modelo}
\label{sec:sd5-5-1}

El modelo separa identidad comercial, lote sanitario, unidad logística, compromiso del pedido y evidencia de entrega. La Figura~\ref{fig:sd5-mr-general} presenta las relaciones generales y las Figuras~\ref{fig:sd5-inventario}--\ref{fig:sd5-dimensional} desarrollan inventario, pedido/preparación, reparto, custodia, cobranza/documentos y analítica. El diccionario 5-A contiene cada atributo; 5-B y 5-C precisan integridad y validación.

\subsection{Dominios de información y almacenes lógicos}
\label{sec:sd5-5-1-1}

El modelo lógico se presenta por dominios: Figuras 5.2--5.6 y A5.1--A5.2 forman sus vistas relacionales; 5.7 presenta el modelo dimensional y A5.3 las estructuras de telemetría, caché y objetos. La Figura 5.1 presenta el MR general de datos de la solución; las Figuras 5.2--5.7 y A5.1--A5.3 desarrollan sus partes con tablas y relaciones específicas. Las vistas por dominio muestran claves y relaciones esenciales; el diccionario 5-A conserva todos los atributos y tipos. Una caja de referencia en otra vista no redefine la tabla ni representa otra base de datos.

Los quince dominios lógicos de S4 distribuyen autoridad y responsabilidad funcional. Maestros y configuración, preventa, rutas, reparto, cobranza/finanzas, Calidad/trazabilidad, EDI, gobierno/acceso y notificaciones residen en los almacenes centrales declarados; inventario y preparación conservan autoridad en PostgreSQL del sitio. Telemetría usa DynamoDB para ingesta térmica y S3/Glue/Redshift para detalle y agregados; BI consulta Redshift/QuickSight. Redis conserva proyecciones y S3\_DOCS evidencia por clase de retención.

El Anexo 5-B integra los quince nombres, módulos, ubicaciones y motores. BD\_RUTAS tiene autoridad central, con itinerarios descargados de consulta; DMS de Talca replica al esquema lector separado. La distinción permite que una operación local siga confirmada en su perímetro sin presentar la proyección central como actualizada durante un corte.

M12 captura telemetría; Calidad decide disposición sanitaria. Una excursión menor transitoria genera advertencia; una crítica o sostenida activa el bloqueo definido en S3/S4 y la resolución final corresponde a Calidad. El ERP conserva emisión tributaria única. Estas responsabilidades evitan tanto dos escritores de un saldo como una resolución sanitaria atribuida al sensor.

\subsection{Modelo conceptual de la operación de Puelche}
\label{sec:sd5-5-1-2}

La Figura 5.1 presenta el MR general con las tablas principales y sus relaciones, agrupadas por dominio; las vistas específicas desarrollan cada parte. Distingue las referencias entre almacenes de las FK físicas y mantiene la vista conceptual de la operación. Las vistas de detalle relacionan productos y lotes con recepción, contenido logístico, custodia, asignaciones del pedido y cumplimiento. El lote se identifica por GTIN canónico+número del proveedor; la unidad tiene UUID interno y SSCC exterior. La composición admite varios lotes por unidad y unidades por lote.

\figurafuente{1.0}{figuras/Fig_5-1_MR_General.png}
  {MR general de datos de la solución}{fig:sd5-mr-general}
  {elaboración propia de LafroX a partir de las decisiones 16.1 N.\textdegree{} 1, N.\textdegree{} 2, N.\textdegree{} 4 y N.\textdegree{} 14 del capítulo 16 del caso (PUCV, 2026c), de los requisitos RT-05.01 y RT-05.23 de las Bases Técnicas Transversales (PUCV, 2026d) y de RF-01.07 y RF-09 del Formulario T-12.}

Cada registro de promesa tiene UUID propio como PK; \texttt{version} ordena sus cambios con unicidad por pedido. La entrega referencia la ORIGINAL del mismo pedido. El pedido conserva ORIGINAL y revisiones como compromisos distintos. Una entrega por pedido admite varios intentos y líneas fragmentadas por asignación/lote/unidad; POD y cobro se vinculan a esos hechos. La cabecera de custodia contiene el evento y su detalle identifica objetos y cantidades. El cliente y sus puntos se desarrollan en 5-M, sin crear una instalación de inventario por cada local de entrega. El documento fiscal emitido por ERP antecede al traslado; el POD acredita el hecho posterior.

\subsection{Maestros compartidos e identidad única}
\label{sec:sd5-5-1-3}

Un maestro tiene dueño funcional, UUID estable e historial; la integración traduce códigos externos mediante equivalencias versionadas y nunca crea identidades paralelas. Producto conserva un GTIN canónico único cuando aplica; otros GTIN/códigos se resuelven por equivalencia aprobada, sin compartir identidad sanitaria entre productos. Corregir descripción o vencimiento no cambia el identificador ni rompe historia.

La Figura A5.1 del Anexo 5-M dibuja maestros, puntos, instalaciones y equivalencias. Un cliente puede tener varios puntos; sitio representa instalación operativa y no cada cliente. Comercial valida productos/clientes, Finanzas condiciones de crédito y Calidad reglas sanitarias; el módulo dueño publica la versión y las copias locales registran instante de actualización. Los conflictos de código se aíslan y concilian antes de publicar, conforme RT-05.09 y 5.2.4.

\subsection{Recepción, inventario y unidad logística}
\label{sec:sd5-5-1-4}

La Figura 5.2 distingue recepción, lote, unidad, contenido, saldo, movimiento, reserva y conteo cíclico. Se exige lote al cierre de las líneas de productos que lo requieren; productos exentos y unidades VACIA/ABIERTA conservan NULL legítimos. El SSCC de 18 dígitos no sustituye el UUID; vencimiento es atributo del lote.

\figurafuente{1.0}{figuras/Fig_5-2_Inventario.png}
  {Recepción, inventario, lote y unidad logística}{fig:sd5-inventario}
  {elaboración propia de LafroX a partir de la decisión 16.1 N.\textdegree{} 2 y N.\textdegree{} 14 del caso, del Subdocumento 4, 4.1.3.6, y de RF-01.07 y RF-02 del Formulario T-12.}

El saldo tiene grano sitio+ubicación+producto+lote+unidad de medida; la conversión se valida antes de comparar o reservar. La transacción bloquea el saldo, valida disponibilidad y registra movimiento/outbox sin duplicar reintentos. Varias asignaciones abastecen una línea de pedido. \texttt{secuencia\_documento} solo existe si el documento de origen aporta ordinal; no se inventa para identificar un hecho histórico. Los ajustes y anulaciones de 5-C conservan causa, dirección y movimiento original; no sobrescriben saldo ni evidencia sin autorización.

\subsection{Preventa, preparación y misión de preparación}
\label{sec:sd5-5-1-5}

En TERRENO, \texttt{preventista\_id} identifica al capturador autorizado; MOSTRADOR, EDI y PORTAL mantienen NULL sin inventar un preventista. La identidad de quien captura o integra se registra en \texttt{gob\_auditoria.actor\_id} con UUID/correlación y \texttt{origen\_captura}: persona autorizada, cuenta de autoatención o actor técnico autenticado, respectivamente. El responsable comercial no se infiere de ese campo. M3 valida cuenta/permisos al confirmar; un pedido pendiente no crea reserva ni promesa de entrega confirmada.

La Figura 5.3 relaciona pedido, líneas, promesas, asignaciones y misiones; la reserva que abastece cada asignación se detalla en la Figura 5.2. La fecha/ventana ORIGINAL se fija al confirmar el pedido y permanece inmutable; una reprogramación agrega REVISION con motivo, actor y autorización. \texttt{fecha\_prometida} y vigencias son timestamps UTC, con zona de origen conocida.

\figurafuente{1.0}{figuras/Fig_5-3_Pedido_Preparacion.png}
  {Pedido, línea, asignación, misión y confirmación de preparación}{fig:sd5-pedido}
  {elaboración propia de LafroX a partir de RNG-01, RNG-15 y RNG-05 del Subdocumento 3, y del Subdocumento 4, 4.1.4.4.}

La promesa conserva las condiciones de S3: 24 h para urbano ingresado antes de las 14:00 y 48 h para rural, periférico o abastecido por cross-docking, según RNG-15. OTIF evalúa ORIGINAL aun después de su cierre de vigencia. La preparación puede dividir una línea entre lotes/unidades y confirmar cantidades parciales; la reserva validada por PostgreSQL del sitio sigue confirmada localmente durante un corte, con consolidación central pendiente. El móvil no reserva contra una copia indicativa de stock.

\subsection{Rutas, reparto, entrega y devolución}
\label{sec:sd5-5-1-6}

La Figura 5.4 separa plan de ruta, viaje, parada, entrega e intento. BD\_RUTAS conserva autoridad central; el itinerario descargado permite ejecutar la ruta autorizada sin señal. Una entrega por pedido admite visitas fallidas, parciales y completas sin sobrescribir las anteriores.

\figurafuente{1.0}{figuras/Fig_5-4_Rutas_Reparto.png}
  {Ruta, viaje, paradas, entrega, intento y devolución}{fig:sd5-rutas}
  {elaboración propia de LafroX a partir de las decisiones 16.1 N.\textdegree{} 1, N.\textdegree{} 3, N.\textdegree{} 10 y N.\textdegree{} 15 del caso, y del Subdocumento 4, 4.1.4.5 y 4.1.4.8.}

Las líneas del intento conservan asignación, lote/unidad, cantidad aceptada/rechazada y causal. La identidad del fragmento se reutiliza en reintentos; distintas asignaciones no se deduplican por compartir producto/cantidad. Una devolución vincula el hecho de origen y su recepción/disposición; la cuenta de envases mantiene saldo por cliente y tipo, sin serialización individual. Los geofences de configuración se distinguen de coordenadas personales GPS, cuya conservación total es doce meses.

\subsection{Trazabilidad sanitaria y cadena de custodia}
\label{sec:sd5-5-1-7}

La Figura 5.5 relaciona evento, objetos, documentos, sensor, asociación temporal y resolución de Calidad. EPCIS 2.0 y CBV 2.0 tienen versiones separadas; \texttt{perfil\_evento} identifica el contrato versionado de la solución descrito en 5-C (GS1, 2022a, 2022b). \texttt{parent\_event\_id} representa causalidad interna y admite NULL en raíz. \texttt{epcis\_parent\_id} identifica el contenedor de una agregación y no un evento causal.

\figurafuente{1.0}{figuras/Fig_5-5_Custodia_EPCIS.png}
  {Cadena de custodia de lote y de unidad logística en formato EPCIS}{fig:sd5-custodia}
  {elaboración propia de LafroX a partir de RF-01.07, RF-09 y RNF-09.02 del Formulario T-12, del Subdocumento 4, 4.1.4.2, y de RT-05.23 según el capítulo 15 del caso.}

Una cabecera puede mencionar varios objetos/cantidades; la asociación temporal vincula sensor, lote, unidad y producto con las lecturas efectivamente aplicables. Regla y calibración se conservan en su versión histórica. Una advertencia menor transitoria no equivale a disposición final; el bloqueo crítico o sostenido se aplica conforme S3/S4 y únicamente Calidad libera o rechaza con fundamento. Lecturas inválidas se preservan marcadas, separando defecto de sensor, excursión y pérdida de cobertura.

La consulta forward sigue lote\,$\rightarrow$\,contenido/unidades \,$\rightarrow$\, objetos/eventos \,$\rightarrow$\, asignaciones \,$\rightarrow$\, intentos/entregas \,$\rightarrow$\, documentos/POD. La backward invierte ese recorrido hasta recepción y proveedor; múltiples recepciones del mismo lote se distinguen por su evidencia, sin atribuir origen único cuando no puede determinarse.

El retiro identifica clientes y documentos en $\le$2~h desde la solicitud de Calidad. Su presupuesto de diseño es 5~min identificación + 10 custodia + 15 entregas/documentos + 10 pendientes + 15 expediente + 15 revisión + 15 margen = 85~min, a demostrar en prueba. Calidad consulta autoridades locales y central; despacho de cada sitio contrasta carga/ruta y conductores aportan pendientes. Los clientes planificados con lote aún offline integran el universo potencial a contactar/bloquear, distinto de entregas confirmadas. A los 30~min sin respuesta se escala a Operaciones; una lista parcial que no identifica ese universo no cumple aceptación. El Anexo 5-J fija evidencia y prueba offline; un sello de última sincronización por sí solo no demuestra completitud.

\subsection{Cobranza, envases, documentos y POD}
\label{sec:sd5-5-1-8}

La Figura 5.6 diferencia cobro, aplicación, autorización, compensación, rendición, documento, carga y POD. El ERP es único emisor fiscal y el documento habilitador precede al traslado; no existe timbre diferido desde el móvil. POD, acuse técnico SII y recepción de mercadería son evidencias distintas, con representación en papel para destinatario no habilitado según S4.

\figurafuente{1.0}{figuras/Fig_5-6_Cobro_Documentos.png}
  {Cobro, rendición, cuenta corriente de envases y documento tributario}{fig:sd5-cobro}
  {elaboración propia de LafroX a partir de RT-16.14 y RT-16.09 de las Bases Técnicas Transversales (PUCV, 2026d), del Anexo 4-U del Subdocumento 4, y de la decisión 16.1 N.\textdegree{} 6 y N.\textdegree{} 12 del caso.}

La identidad fiscal usa tipo/emisor/folio/versión; hash verifica contenido y las referencias tipadas comprueban la entidad de negocio. POD admite varias evidencias por entrega/intento; el móvil puede cerrar localmente bajo autorización con evidencia durable, conservando estado de verificación central posterior.

Un timeout de POS mantiene pago incierto: consulta, conciliación y eventual cancelación/compensación dependen del contrato del tercero. Un pago alternativo necesita autorización financiera y vínculo al primero; agotar reintentos no prueba rechazo ni evita por sí mismo doble cargo. Autoriza un actor con rol financiero distinto del ejecutor. \texttt{ajuste\_redondeo} expresa importe registrado menos importe externo comprobado solo para tarjeta; efectivo/crédito no fabrican autorización externa. Las cuentas de envases son saldos por cliente/tipo. S4, Anexo 4-M, define AL-DTE-01 y AL-DR-01; su aceptación requiere ensayo.

\subsection{Intercambio electrónico, notificaciones y gobierno del acceso}
\label{sec:sd5-5-1-9}

La Figura A5.2 de 5-M desarrolla perfil de cadena, equivalencias, mensajes/ASN, avisos, asignación de roles, auditoría y excepciones. El perfil establece esquema y campos admitidos; se conserva mensaje original, hash, identidad externa y correlación. Un reintento devuelve resultado previo sin duplicar efecto. Un código de cadena se traduce a identidad canónica; no crea cliente/producto por defecto.

La auditoría registra actor, dispositivo, instante, operación, resultado y valores anteriores/posteriores mínimos. La excepción conserva contexto histórico, regla y versión; su identidad es (\texttt{correlacion\_id}, \texttt{clave}), no hora de llegada. Los mensajes/notificaciones usan solo integraciones habilitadas en S4 y su estado de envío no sustituye recepción de mercadería ni acuse fiscal.

\subsection{Telemetría, caché y objetos}
\label{sec:sd5-5-1-10}

La Figura A5.3 del Anexo 5-M distingue ingesta térmica cruda en DynamoDB, archivo detallado S3 Parquet/Glue, serie/analítica Redshift, posición personal y objetos. Raw dura treinta días; el perfil detallado y las versiones que lo interpretan permanecen cinco años. El agregado no permite reconstruir una lectura perdida, por lo que no sustituye el detalle.

S4, Anexo 4-W.4, fija 21 puntos térmicos $\times$288 y 28 termógrafos $\times$162 = 10.584 muestras/día. Operan tres gateways: dos en Talca y uno en Concepción; las reservas de T-11 no agregan ingesta ordinaria. El sensor referencia su gateway, sin crear otra autoridad de maestros.

GPS permanece en sa-east-1, raw treinta días y plazo personal total doce meses, sin réplica internacional. La posición no se confunde con geocerca de configuración. S3\_DOCS conserva originales por clase; Redis contiene lectura reconstruible y Room/SQLite conserva además hechos/evidencias durables. Las claves, vigencias y purgas están en 5-G.

\subsection{Modelo lógico del dominio analítico}
\label{sec:sd5-5-1-11}

La Figura 5.7 define hechos y dimensiones para analizar servicio, inventario, cobro y frío. El grano e identidad de origen evitan duplicados de carga: pedido para compromiso, entrega/intento para ejecución y línea/asignación para cantidades; la entrega por pedido no cuenta cada visita como pedido independiente.

\figurafuente{1.0}{figuras/Fig_5-7_Dimensional.png}
  {Modelo dimensional de explotación}{fig:sd5-dimensional}
  {elaboración propia de LafroX a partir de RF-11.06 y RF-11.07 del Formulario T-12, de los criterios R18-04 y R18-11 del capítulo 18 del caso, y de RNG-09 del Subdocumento 3.}

El modelo incluye dimensiones de fecha, cliente, producto, sitio, ruta y canal; cliente, producto, sitio y ruta conservan versiones y vigencias. El cálculo previo al grano del pedido evita multiplicar OTIF/importes al unir líneas; cada hecho conserva origen y controles monetarios/GPS. Redshift/QuickSight sirve analítica pesada y la API consultas operativas. 5-H fija fórmulas, audiencia, filtros, profundización y latencia; costo de servir inicia E2. RT-05.10/.24/.30 deseables conservan su alcance no ofertado.

\subsection{Reglas relacionales resueltas}
\label{sec:sd5-5-1-12}

Los extremos opcionales representan estados iniciales o excepciones, no omisión de controles: TERRENO exige preventista; producto con lote exigido exige lote; raíz causal es el único evento sin padre; un sensor fijo exige sitio y uno móvil vehículo. La unidad VACIA/ABIERTA puede no tener contenido; al cerrarse exige contenido y SSCC según el hito. Una misión puede no tener confirmación antes del picking, pero no cierra sin validar sus líneas. Pedido PENDIENTE no tiene promesa confirmada; al confirmarse mantiene exactamente una ORIGINAL y las revisiones autorizadas. Cada pedido admite cero o una entrega y varios intentos vinculados a esa entrega.

El modelo relacional lógico se presenta por dominios en las Figuras 5.2--5.6, con atributos, claves y cardinalidades; la Figura 5.7 desarrolla el modelo dimensional analítico. Las Figuras A5.1--A5.3 de 5-M completan maestros, EDI/gobierno y telemetría, y 5-A contiene el diccionario exacto. Las FK físicas se limitan al mismo almacén; entre dominios, sedes, motores y objetos se validan identidad/versión y conciliación.

Lote usa GTIN canónico+número, con vencimiento como atributo. Unidad usa UUID y SSCC exterior, composición de varios lotes y estados VACIA/ABIERTA permitidos. Saldo tiene grano sitio+ubicación+producto+lote+unidad; recepción conserva \texttt{secuencia\_documento} solo si el origen entrega ordinal. Reserva y asignación remiten a la línea de pedido concreta y permiten fragmentación por lote/unidad.

Custodia separa cabecera y objetos; \texttt{parent\_event\_id} representa genealogía causal y admite NULL en raíz, mientras la agregación logística representa contenedor/contenido. \texttt{ts\_evento} y \texttt{secuencia\_linea} del detalle mantienen timestamp de cabecera y ordinal estable; las restricciones compuestas de 5-F hacen implementable la partición sin duplicar un reintento entre meses.

Un pedido tiene una entrega con varios intentos; las líneas del intento distinguen asignación, lote/unidad y cantidades aceptadas/rechazadas. Promesa ORIGINAL es inmutable y la fecha/ventana usa UTC; la entrega referencia \texttt{promesa\_id}, sin columna duplicada \texttt{fecha\_promesa\_original}. Las revisiones autorizadas no reemplazan la base de OTIF.

Documento fiscal, POD, acuse técnico y recepción de mercadería conservan identidad propia. La clave fiscal incluye tipo/emisor/folio/versión; el hash verifica contenido. El ERP emite antes del traslado. Cobro y compensación se concilian con segregación de autorizador, nunca mediante una corrección silenciosa de importe.

El Anexo 5-C define parámetros tipados, perfiles EPCIS/CBV separados, movimientos/reversiones y vigencias sin solapamiento; 5-F declara índices y unicidad global. Los casos de borde se comprobarán en PREPROD según 5-J, sin afirmar que el diseño documental ya equivale a restricciones ejecutadas.
